\newpage
% \color{white}
% \pagecolor{black}

\ifdefined\useChapters
\chapter{Sacrifícios}
\else
\chapter{}
\fi

— Pra onde a gente vai?

— Você fala demais. Só me segue — disse Ítalo, volitando tão rápido que o vento arrancou a resposta da minha boca.

Ítalo voava com a naturalidade de quem nasceu no ar.  
Eu, apesar de ter aprendido muito nos últimos dias, ainda me sentia como alguém que tenta correr em um chão que muda de lugar.  
Precisava de concentração demais para manter o ritmo — e concentração demais significava menos tempo pra pensar no que estávamos prestes a fazer.

E, pra ser sincero, pensar era a última coisa que eu queria.

— Eu não gosto dessa ideia de matar alguém — murmurei, tentando não perder altitude.

Ítalo desacelerou por um segundo, o suficiente para que sua voz voltasse até mim.

— Então aprende a gostar de viver. Às vezes é isso ou ser destruído — disse, e acelerou de novo.

A cidade desapareceu sob nossos pés, dando lugar ao vilarejo distante.  
Senti o estômago afundar ao reconhecer o caminho, os telhados queimados, as sombras.

E, logo abaixo, a casa onde tudo começou.

Ítalo apontou para ela.

— Ali. É onde a gente precisa estar.

— Você só pode estar louco — respondi, meu coração ralando no peito. — Deve estar todo mundo lá. A gente vai virar pó em segundos.

— Conheço a rotina deles. — Ele falou como quem recita a hora exata em que o sol nasce. — Crispim nunca muda nada. Hoje, só tem lá quem a gente precisa pegar.

E, como se o mundo estivesse sob seu comando, Ítalo desceu em queda livre.

Tentei acompanhar.  
Falhei.  
Quase bati no telhado antes de reencontrar o ar.

Ele entrou pela claraboia — a mesma por onde fugi da primeira vez.  
Eu segui atrás, sabendo que, se desse errado, talvez fosse minha última vez ali.

Quando caí na sala, o ar me fugiu pelos pulmões.  
Ítalo já estava com Joana imobilizada, uma mão firme no braço dela, a outra ajustando a faca no pescoço.

Tudo tinha acontecido rápido demais. Fácil demais.

— Você não acha que está tudo fácil demais? — sussurrei.

Ítalo deu um sorriso curto.

— Crispim é um ogro. Nunca leu um livro. Pra gente isso é clichê. Pra ele, é novidade.

Joana tremia.  
Eu não sabia se era raiva, medo ou os dois.

— Não machuca ela — pedi. — Ela não tem nada a ver com isso.

— Como não tem? — Ítalo rebateu, sem tirar a faca da pele dela. — Ela sabia de tudo. Sempre soube.

As palavras dele caíram como marteladas.

— É mentira! — Joana gritou, a voz quebrada. — Eu não sabia de nada. Meu pai me enganou. Eu juro.

Ela olhou nos meus olhos.  
E havia uma dor ali que parecia verdadeira — uma mistura de vergonha e desespero.

Antes que eu respondesse, uma voz explodiu no corredor.

— O que é isso?! — Crispim surgiu correndo, o rosto deformado entre espanto e horror.

Ítalo puxou Joana com força.

— Você não devia estar aqui, seu velho escroto.

Uma faca reluziu. A lâmina tocou a pele dela.

E então—

Palmas.  
Lentas.  
Cruéis.

O garçom entrou.

Sempre ele.  
Sempre invisível até decidir aparecer.

— Vocês são incríveis — disse, com um sorriso torto. — Acham mesmo que podem impedir Crispim de brincar com a cabeça de vocês?

Ele ergueu o telefone de Joana entre dois dedos.

— Quando ouvi a ligação que ela fez pra você, imaginei que tentariam alguma coisa — continuou. — Vocês chegaram mais tarde do que eu esperava. Fiquei escondido no corredor esse tempo todo, vendo o plano perfeito de vocês entrar pela claraboia.

Crispim não se moveu.  
Estava paralisado — não pelo garçom, mas pelo medo de perder a filha.

— Vai, Crispim — o garçom provocou. — Mostra o que você sabe fazer.

Nada.

— Anda — repetiu ele, aproximando-se. — Ou você é só conversa?

O silêncio do Crispim doeu.  
Mais que qualquer grito.

Eu tentei acender chamas.  
Nada aconteceu.  
Meu corpo estava mudo, como se minhas mãos fossem de outra pessoa.

O garçom olhou pra mim e sorriu.

— Pronto. Agora sim. Isso é o que eu gosto de ver.

Ele tinha tomado minhas habilidades.  
Como se fosse fácil.  
Como se fosse natural.

Crispim deu um passo para trás, completamente vulnerável.

— Você é um frouxo — disse o garçom, sem olhar para ele. — Sempre foi.

Joana começou a chorar.  
Ítalo pressionou a faca, mas não o suficiente para cortar.

O garçom suspirou, como quem faz um favor pro universo.

— Para conseguir o que queremos, alguns sacrifícios precisam ser feitos.  
E já que você não tem coragem…

Tudo aconteceu em um segundo que pareceu uma eternidade:

Um brilho metálico.  
Um movimento suave.  
O gatilho apertado sem hesitação.

O tiro.  
O impacto seco no peito de Joana.  
Um som abafado, como um suspiro que não conseguiu sair.

O corpo dela caiu como se todas as cores do mundo tivessem sido drenadas de uma vez.

Crispim gritou — um grito que não parecia humano, de tão rasgado.

Ele correu, abraçando a filha no chão, tremendo como se quisesse voltar no tempo pela força das mãos.

Eu senti o estômago virar.

Ítalo me puxou pelo braço.

— Vamos sair daqui. Agora. Sem habilidades e com ele armado, a gente tá morto.

— Eu não posso mais fugir — respondi, a voz quebrada.

— Pode sim. Tem que poder. Vamos!

Mas algo dentro de mim já tinha cruzado uma linha.  
Não era coragem.  
Era algo mais antigo.  
Um tipo de certeza que não vem do pensamento — vem do limite.

Do topo da sala, o garçom nos observava.

— Vai lá, herói — disse, provocando. — Mostra do que é feito.

Eu tremia.  
Mas não recuei.

Minha mão subiu, não por força, mas por instinto.

A arma saiu da mão do garçom, arrancada pelo ar, e ficou suspensa entre nós dois, flutuando, apontada diretamente para a testa dele.

— O que é isso, Crispim?! — o garçom gritou, sem tirar os olhos da arma. — Você disse que ele só criava ilusões!

Crispim não respondeu.  
Só se encolheu ainda mais sobre o corpo da filha.

— E então? — provocou o garçom, empurrando a própria testa contra o cano. — Vai me matar?  

— Acha mesmo que é um herói?

Eu não conseguia puxar o gatilho.  
Ele sabia disso.  
Eu sabia disso.

— Você não tem coragem — ele disse, firme como uma sentença.

Uma mão pousou no meu ombro.  
O chão girou.  
O ar sumiu.  
A visão estourou em branco.

E então — nada.

A arma caiu no tapete.

E a sala ficou só com três figuras:

Crispim, ajoelhado.  
O corpo morto de Joana.  
E o garçom — vivo, imóvel e sorrindo.
